% UNIFIED 2026-05-13
% RE-AUDIT WARNING — DATE 2026-05-14
% The endpoint repair below should not be read as an unconditional closure.
% Two issues remain after re-audit:
%   (1) the small-window trace argument gives an O(sqrt(delta_T)) metric
%       Markov contribution, not O(delta_T);
%   (2) Lemma `lem:wmt:A0` uses a perimeter-free dilation Lipschitz estimate
%       that is false for rough bounded measurable sets.
% The corrected route-status digest reduces the remaining obstruction to the
% bad-(-t') Lebesgue kinetic estimate
%   int_{B_tau} |dot F_rho|_X^2 d rho <= C delta_T.
% Keep this subsection as a draft until that estimate is supplied or the
% trace transport step is replaced.
% AUDIT B1/M1/M2 REPAIR — DATE 2026-05-13
% -----------------------------------------------------------------------------
% This file has been edited to address the BLOCKING (B1) and MAJOR (M1, M2)
% findings of `Manuscript/audit_plan2.md` (Agent 16). Summary of changes:
%   B1: Theorem `thm:wmt:trace` now instantiates the abstract weighted-trace
%       lemma on the window `[rho_delta - c_0 sqrt(delta_T), rho_delta]` with
%       L* = c_0 sqrt(delta_T), placing widehat-rho within an
%       O(sqrt(delta_T))-neighbourhood of rho_delta. (H2) is verified on the
%       new window using a smallness threshold delta_T <= delta_0(n,R,rho_*).
%       Closing remark is rewritten so that |Omega \ E_{widehat-rho}| =
%       O(sqrt(delta_T)).
%   M1: Lemma `lem:fj-center-package` now explicitly identifies z_rho := bar
%       z_rho (bulk centroid), citing Agent 11 Remark `rmk:centre-selection`.
%   M2: kappa = kappa(n, R, rho_*) is pinned canonically (per notation
%       register §12). All constants C_*^act, C_*^kin, C_* now have the
%       signature C_*(n, R, rho_*).
% See `Manuscript/audit_plan2_repair_report.md`.
% -----------------------------------------------------------------------------
% SOURCE MAP (Agent 12, Plan 2 §4.4 — Integrated Action and Weighted Metric Trace)
% -----------------------------------------------------------------------------
% Local building blocks consumed (LOCAL):
%   - Plan 2/WIP/WIP_WeightedMetricTrace.tex
%       * Theorem (integrated action bound, eq:52int)
%       * Theorem (kinetic bound K, eq:K)
%       * Theorem (weighted metric trace, eq:trace)
%       * Lemmas: FJ centre package, radial-slicing, good-level homothetic
%         trace, uniform metric Lipschitz (A0), bad-set lemma
%   - Plan 2/WIP/WIP_MetricFramework.tex
%       * Theorem T: |F'_rho|_X <= rho_*^{-n} inf_a int |V_rho - H_{z,rho}
%                       - a . nu_rho| dH^{n-1}
%       * AGS absolute continuity / metric derivative (Thm 1.1.2)
%   - Plan 2/WIP/WIP_LevelSetIdentity.tex
%       * Exact deficit identity: (2/|Omega|) (E(Omega) - E(B^*))
%           = (1/|Omega|) int (D_H + D_I)/(c_n m(t)^{1-2/n}) dt
%       * Profile-gap corollary used for (H2) and bad-set
%   - Plan 2/agent4-weighted-metric-trace
%       * Abstract weighted metric trace lemma (proof of Thm in §4)
%       * Bad-set Lemma 3.1; Corollary 3.2
%   - Plan 2/wave2-B-kinetic-bound
%       * Uniform L1-Lipschitz bound (A0), Lemma 2.1
%       * Bad-(-t') Lemma 4.1; integrated Talenti gap (3.3)
%   - Plan 2/wave3-G-route-delta-assembly.md  (verification that the
%       integrated chain consumes (R) only against d\mu, no rate loss)
%
% Cited from the literature (CITE; see Manuscript/04_references_audit.md):
%   - Ambrosio-Gigli-Savaré 2008, Thm 1.1.2          [AGS2008]
%   - Fusco-Julin 2014, Calc. Var. PDE 50            [FJ2014]
%   - Talenti 1976, Ann. Mat. Pura Appl. (4) 110     [Talenti1976]
%   - Maggi 2012 (sets of finite perimeter)          [Maggi2012]
%   - Ambrosio-Fusco-Pallara 2000 (BV book)          [AFP2000]
%   - Figalli-Maggi 2011, Appendix A, Thm A.1        [FigalliMaggi2011]
%
% Imported by Agent 19 into Manuscript/03_plan2_draft.tex as Section §4.4.
% Constants tracked: every C, c is annotated with its dependence on (n, R, rho_*).
% (Per the M2 repair dated 2026-05-13, kappa = kappa(n, R, rho_*) is canonical,
% so all C_*^{act}, C_*^{kin}, C_* have signature (n, R, rho_*).)
% -----------------------------------------------------------------------------

\subsection{Integrated action and the weighted metric trace}\label{subsec:weighted-metric-trace}

This subsection chains the level-set deficit identity of
Section~\ref{subsec:level-set-identity}, the metric first-variation
estimate of Section~\ref{subsec:metric-framework}, the scaled
Fusco--Julin inequality and the oriented radial trace lemma of
Section~\ref{subsec:fj-radial-trace} into the integrated action bound
\eqref{eq:wmt:52int} and, via a good/bad isoperimetric-defect
decomposition, into the Lebesgue kinetic bound \eqref{eq:wmt:K}. The
final output is the \emph{weighted metric trace}: a radius
$\rhad\in[\rstar,\rhodelta)$ at which the rescaled level set
$F_\rhad=\rhad^{-1}E_\rhad$ is metrically close to the unit ball at the
sharp $\deltaT$-rate; see Theorem~\ref{thm:wmt:trace}.

Throughout this subsection we adopt the standing hypotheses of
Section~\ref{sec:plan2}: $n\ge 2$, $R\ge 2$, $\rstar\in(0,1)$, an open set
$\Om\subset B_R$ with $|\Om|=\omega_n$, variational torsion function
$u=u_\Om$, $\rho$-foliation $E_\rho=\{u>t(\rho)\}$ with
$|E_\rho|=\omega_n\rho^n$ for $\rho\in[\rstar,1]$, rescaled level sets
$F_\rho:=\rho^{-1}E_\rho\subset B_{R/\rstar}$, and the measure
$\mu(d\rho):=(-t'(\rho))\,d\rho$ on $[\rstar,1]$. We write
$\delta:=\deltaT(\Om)$.

\paragraph{Canonical choice of $\kappa$.}
Per the notation register \S12, the boundary-layer constant $\kappa$ is
a fixed dimensional choice $\kappa=\kappa(n,R,\rstar)>0$. Throughout
\S4.4 we pin $\kappa$ to this canonical value; explicitly, we set
\[
   \kappa=\kappa(n,R,\rstar):=1
\]
(any fixed positive constant of size $\Theta_n(1)$ works; the value
$\kappa=1$ is the simplest choice consistent with the notation
register). All subsequent constants $C_*^{\mathrm{act}},
C_*^{\mathrm{kin}}, C_*$ then become functions of $(n,R,\rstar)$ alone;
the smallness threshold $\delta_0$ likewise depends only on
$(n,R,\rstar)$. We define
\[
   \rhodelta:=1-\kappa\sqrt{\delta}.
\]
Throughout we assume
\begin{equation}\label{eq:wmt:delta-small}
   \delta\le\delta_0(n,R,\rstar),
\end{equation}
with $\delta_0(n,R,\rstar)>0$ chosen so that
$\rhodelta\ge\rstar+\tfrac12(1-\rstar)$ and so that the additional
smallness requirements (B1) of Theorem~\ref{thm:wmt:trace} below hold;
for $\delta$ outside this regime the conclusion is vacuous or follows
trivially from the FMP fallback $\Fr(\Om)^2\le C(n)\deltaT(\Om)$ for
$\deltaT(\Om)\ge\delta_0$.

The level-set deficit identity of Section~\ref{subsec:level-set-identity}
reads, in our $\rho$-parametrisation,
\begin{equation}\label{eq:wmt:DH-DI-int}
   \int_{\rstar}^{1}
      \bigl(\mathcal H(\rho)+\mathcal I(\rho)\bigr)\,d\mu(\rho)
   \;\le\; C_1(n,\rstar)\,\delta,
\end{equation}
where we have written
\[
   \mathcal H(\rho):=D_H(t(\rho)),\qquad
   \mathcal I(\rho):=D_I(t(\rho)),
\]
and used $dt=-t'(\rho)\,d\rho=d\mu$. The explicit constant is
$C_1(n,\rstar)=2\,c_n\,(\omega_n\rstar^n)^{1-2/n}/\omega_n
=2n^2\omega_n^{2/n}(\omega_n\rstar^n)^{1-2/n}/\omega_n$; see
Theorem~\ref{thm:lsid:identity}.

Since $P(E_\rho)+P(B_\rho)\ge 2P(B_\rho)=2n\omega_n\rho^{n-1}\ge
2n\omega_n\rstar^{n-1}$ and $D_I(t(\rho))
=(P(E_\rho)-P(B_\rho))(P(E_\rho)+P(B_\rho))$, the perimeter defect
satisfies
\begin{equation}\label{eq:wmt:DP-DI}
   0\le\Delta P_\rho
        :=P(E_\rho)-P(B_\rho)
   \le \frac{\mathcal I(\rho)}{2n\omega_n\rstar^{n-1}}
   =: C_2(n,\rstar)\,\mathcal I(\rho).
\end{equation}

\paragraph{Notation for this subsection}

\begin{itemize}
\item $\gamma(\rho):=[F_\rho]\in\X$ is the rescaled curve in the metric
      quotient $(\X,\dX)$ of Section~\ref{subsec:metric-framework}.
      AGS \cite[Thm.~1.1.2]{AGS2008} provides the metric derivative
      $|\dot\gamma|(\rho)=|\dot F_\rho|_{\X}$.
\item $z_\rho$ denotes the bulk centroid of $E_\rho$, identified in
      Lemma~\ref{lem:fj-center-package} below with the canonical centre
      $\zbar=\bar z_\rho$ of the notation register \S10. Outside the
      good set $G_I$ defined in \eqref{eq:wmt:def-GI} we set
      $z_\rho=0$.
\item $\bar V_\rho:=P(E_\rho)^{-1}\int_{\partial^*E_\rho}V_\rho\,
      d\mathcal H^{n-1}=P(B_\rho)/P(E_\rho)$, which equals
      the mean value of $V_\rho$ on $\partial^*E_\rho$ and satisfies
      $1-\bar V_\rho=\Delta P_\rho/P(E_\rho)$.
\item $H_{z_\rho,\rho}(x):=(x-z_\rho)\cdot\nu_\rho(x)/\rho$ is the
      homothetic radial field of Section~\ref{subsec:metric-framework}.
\end{itemize}

\subsubsection{4.4.1 Good and bad isoperimetric-defect levels}
\label{subsec:wmt:good-bad}

\begin{definition}[Good and bad isoperimetric-defect levels]
\label{def:wmt:GI-BI}
Let $\eta_I=\eta_I(n,R,\rstar)>0$ be the Fusco--Julin smallness
threshold from Lemma~\ref{lem:fj-center-package} below. Define
\begin{align}
   G_I &:= \{\rho\in[\rstar,\rhodelta]:\mathcal I(\rho)\le\eta_I\},
       \label{eq:wmt:def-GI}\\
   B_I &:= [\rstar,\rhodelta]\setminus G_I
       =\{\rho\in[\rstar,\rhodelta]:\mathcal I(\rho)>\eta_I\}.
       \label{eq:wmt:def-BI}
\end{align}
\end{definition}

The threshold $\eta_I=\eta_I(n,R,\rstar)$ will be fixed at the end of
Lemma~\ref{lem:fj-center-package}: it is the smallest of the
Fusco--Julin smallness scale on $\{|E|=|B_\rho|,\rho\in[\rstar,1]\}$ (see
Proposition~\ref{prop:fj-scaled}) and the elementary smallness
threshold needed for the FJ centre to lie in $B_{2R}$. In particular,
$\eta_I$ does \emph{not} depend on $\delta$ or on $\Om$.

\begin{lemma}[Markov bound on $B_I$]\label{lem:wmt:markov-BI}
Under \eqref{eq:wmt:delta-small},
\begin{equation}\label{eq:wmt:markov-BI}
   \mu(B_I)
   \;\le\; \eta_I^{-1}\int_{B_I}\mathcal I(\rho)\,d\mu(\rho)
   \;\le\; \eta_I^{-1}\int_{\rstar}^{1}
              \bigl(\mathcal H+\mathcal I\bigr)\,d\mu
   \;\le\; \eta_I^{-1}\,C_1(n,\rstar)\,\delta.
\end{equation}
\end{lemma}

\begin{proof}
By Definition~\ref{def:wmt:GI-BI}, $\mathcal I(\rho)>\eta_I$ on $B_I$,
hence
$\mu(B_I)=\int_{B_I}d\mu \le \eta_I^{-1}\int_{B_I}\mathcal I\,d\mu$.
Both $\mathcal H,\mathcal I\ge 0$, so the second inequality is trivial.
The last inequality is \eqref{eq:wmt:DH-DI-int}.
\end{proof}

\begin{remark}[The sharp rate is preserved]\label{rem:wmt:rate-good-bad}
The Markov bound \eqref{eq:wmt:markov-BI} is the only place where the
size of $B_I$ enters; the rate is $\mu(B_I)=O(\delta)$, not
$O(\sqrt{\delta})$. The good-set quantitative estimates
\eqref{eq:wmt:FJ-asym}--\eqref{eq:wmt:FJ-normal} below are pointwise in
$\mathcal I(\rho)$, which is then integrated against $d\mu$ via
\eqref{eq:wmt:DH-DI-int}. No square root of $\mathcal I$ is integrated
unweighted; consequently the integrated bound \eqref{eq:wmt:52int}
preserves the sharp $\delta$ rate. This is the point at which the
naive route would lose a $\sqrt{\delta}$.
\end{remark}

\begin{lemma}[FJ centre package, fixed to the bulk centroid]
\label{lem:fj-center-package}\label{lem:FJ-package}
There exist $\eta_I=\eta_I(n,R,\rstar)>0$ and
$C_{\mathrm{FJ}}=C_{\mathrm{FJ}}(n,R,\rstar)$ such that, defining the
centre field on $G_I$ as the bulk centroid
\begin{equation}\label{eq:wmt:zrho-is-centroid}
   z_\rho:=\zbar=\bar z_\rho=|E_\rho|^{-1}\int_{E_\rho}x\,dx,
   \qquad \rho\in G_I,
\end{equation}
and $z_\rho:=0$ on $B_I$, the map $\rho\mapsto z_\rho$ is Borel
measurable on $[\rstar,\rhodelta]$, satisfies $|z_\rho|\le 2R$ on
$G_I$, and for every $\rho\in G_I$,
\begin{align}
   |E_\rho\Delta B_\rho(z_\rho)|^2
      &\le C_{\mathrm{FJ}}\,\mathcal I(\rho),\label{eq:wmt:FJ-asym}\\
   \int_{\partial^*E_\rho}
      \Bigl|\nu_\rho(x)-\frac{x-z_\rho}{|x-z_\rho|}\Bigr|^2
      d\mathcal H^{n-1}(x)
      &\le C_{\mathrm{FJ}}\,\mathcal I(\rho).\label{eq:wmt:FJ-normal}
\end{align}
Moreover the centroid containment
\(B(z_\rho,\rho/4)\subset E_\rho\subset B(z_\rho,3\rho)\) of
Lemma~\ref{lem:centroid-borel} (Section~\ref{subsec:fj-radial-trace})
holds on $G_I$, justifying the use of the oriented radial trace lemma
at $z_\rho$.
\end{lemma}

\begin{proof}
The choice $z_\rho:=\zbar$ on $G_I$ is the canonical one of the
notation register \S10, cf.\ also Agent~11
Remark~\ref{rmk:centre-selection} of
Section~\ref{subsec:fj-radial-trace}. By that remark, the centroid
$\zbar$ is admissible in the Fusco--Julin bound \eqref{eq:plan2-FJ-scaled}
up to absorbing the centroid--Fraenkel offset
$|\zbar-z^{\mathrm{Fr}}_\rho|\le C(n,\rstar)\sqrt{\delta_T}$ into the
constant; the smallness threshold \eqref{eq:wmt:delta-small} (with
$\delta_0$ shrunk if necessary so that
$C(n,\rstar)\sqrt{\delta_0}\le\rho/4$) is precisely what makes this
absorption legitimate. The Fraenkel bound
\eqref{eq:wmt:FJ-asym} and the normal-oscillation bound
\eqref{eq:wmt:FJ-normal} at $z_\rho=\zbar$ are then
Lemmas~\ref{lem:fraenkel-centroid} and~\ref{lem:normal-osc-centroid}
of Section~\ref{subsec:fj-radial-trace}, combined with the divergence
identity \(\int_{\partial^*E_\rho}|\nu_\rho-(x-z_\rho)/|x-z_\rho||^2\,
d\mathcal H^{n-1}=2\beta(E_\rho,z_\rho)^2\) (see
Theorem~\ref{thm:fj-radial-divergence}). The Lipschitz/Borel
regularity of $\rho\mapsto\zbar$ is Lemma~\ref{lem:centroid-borel};
$|\zbar|\le 2R$ follows from $\zbar\in E_\rho\subset B_R$. The
threshold $\eta_I=\eta_I(n,R,\rstar)$ is the smaller of the
Fusco--Julin smallness scale on $\{|E|=|B_\rho|,\rho\in[\rstar,1]\}$
and the threshold ensuring centroid containment
\(B(\zbar,\rho/4)\subset E_\rho\subset B(\zbar,3\rho)\) of
Lemma~\ref{lem:centroid-borel}.
\end{proof}

\begin{lemma}[Oriented radial trace at the FJ centre]
\label{lem:wmt:radial-trace}
For every $\rho\in G_I$,
\begin{equation}\label{eq:wmt:radial-trace}
   T_\rho:=\int_{\partial^*E_\rho}
      \Bigl|\frac{|x-z_\rho|}{\rho}-1\Bigr|\,d\mathcal H^{n-1}(x)
   \;\le\; C_3(n,R,\rstar)\,\mathcal I(\rho)^{1/2}.
\end{equation}
\end{lemma}

\begin{proof}
This is Lemma~\ref{lem:radial-trace-master} of
Section~\ref{subsec:fj-radial-trace}, applied with the centre field of
Lemma~\ref{lem:fj-center-package} and combined with
\eqref{eq:wmt:DP-DI}, \eqref{eq:wmt:FJ-asym}, and
\eqref{eq:wmt:FJ-normal}. Since $\mathcal I(\rho)\le\eta_I\le 1$ on
$G_I$, the FJ asymmetry square-root piece $\mathcal I(\rho)^{1/2}$
dominates the linear piece $\mathcal I(\rho)$.
\end{proof}

\begin{lemma}[Good-level homothetic trace]\label{lem:wmt:good-htrace}
For every $\rho\in G_I$,
\begin{equation}\label{eq:wmt:htrace-good}
   \Biggl(\int_{\partial^*E_\rho}
       |H_{z_\rho,\rho}(x)-\bar V_\rho|\,
       d\mathcal H^{n-1}(x)\Biggr)^2
   \;\le\; C_4(n,R,\rstar)\,\mathcal I(\rho).
\end{equation}
\end{lemma}

\begin{proof}
Set $e_r(x):=(x-z_\rho)/|x-z_\rho|$. Then
$H_{z_\rho,\rho}(x)
   =(x-z_\rho)\cdot\nu_\rho(x)/\rho
   =(|x-z_\rho|/\rho)\,e_r\cdot\nu_\rho$.
A direct algebraic identity gives
\[
   H_{z_\rho,\rho}-1
   =\Bigl(\frac{|x-z_\rho|}{\rho}-1\Bigr)\,e_r\cdot\nu_\rho
     +(e_r\cdot\nu_\rho-1).
\]
Since $|e_r|=|\nu_\rho|=1$, $|e_r\cdot\nu_\rho-1|
=\tfrac12|e_r-\nu_\rho|^2$. Hence
\[
   \int_{\partial^*E_\rho}|H_{z_\rho,\rho}-1|\,d\mathcal H^{n-1}
   \le T_\rho+\tfrac12\int_{\partial^*E_\rho}|e_r-\nu_\rho|^2\,
       d\mathcal H^{n-1}.
\]
By Lemma~\ref{lem:wmt:radial-trace} and \eqref{eq:wmt:FJ-normal},
\(
   \int_{\partial^*E_\rho}|H_{z_\rho,\rho}-1|\,d\mathcal H^{n-1}
   \le C_3\,\mathcal I^{1/2}+\tfrac12 C_{\mathrm{FJ}}\,\mathcal I
   \le C_5(n,R,\rstar)\,\mathcal I^{1/2}.
\)
On the other hand
\[
   P(E_\rho)|1-\bar V_\rho|
   =P(E_\rho)-P(B_\rho)
   =\Delta P_\rho
   \le C_2\,\mathcal I(\rho)
   \le C_2\,\mathcal I(\rho)^{1/2},
\]
using $\mathcal I\le\eta_I\le 1$ on $G_I$. Triangle inequality gives
$\int|H_{z_\rho,\rho}-\bar V_\rho|\,d\mathcal H^{n-1}
\le C_6(n,R,\rstar)\,\mathcal I(\rho)^{1/2}$, which squares to
\eqref{eq:wmt:htrace-good} with
$C_4(n,R,\rstar):=C_6(n,R,\rstar)^2$.
\end{proof}

\subsubsection{4.4.2 The integrated action bound}
\label{subsec:wmt:integrated-action}

We now assemble \eqref{eq:wmt:DH-DI-int},
Lemma~\ref{lem:fj-center-package}, Theorem~T of
Section~\ref{subsec:metric-framework}, and the good-level homothetic
trace into the integrated action bound.

\begin{theorem}[Integrated action bound]\label{thm:wmt:52int}
There is $C_*^{\mathrm{act}}=C_*^{\mathrm{act}}(n,R,\rstar)$
such that
\begin{equation}\label{eq:wmt:52int}
   \int_{\rstar}^{\rhodelta}
     \Bigl(\dX(F_\rho,B_1)^2+|\dot F_\rho|_{\X}^2\Bigr)\,d\mu(\rho)
   \;\le\; C_*^{\mathrm{act}}\,\deltaT(\Om).
\end{equation}
\end{theorem}

\begin{proof}
We bound the two integrands separately. For both, we estimate
pointwise on $G_I$ and use the Markov bound on $B_I$.

\emph{Distance term.} On $G_I$, the quotient pseudodistance satisfies
\begin{equation}\label{eq:wmt:dist-trivial}
   \dX(F_\rho,B_1)
   =\dX([\rho^{-1}E_\rho],[B_1])
   \le \rho^{-n}\,|E_\rho\Delta B_\rho(z_\rho)|
   \le \rstar^{-n}\,|E_\rho\Delta B_\rho(z_\rho)|,
\end{equation}
because $\rho^{-n}\,\mathbf{1}_{\rho^{-1}A}=\rho^{-n}\mathbf{1}_A\circ
(\rho\,\cdot\,)$ and the Lebesgue measure of $\rho^{-1}A$ equals
$\rho^{-n}|A|$ for any measurable $A$ (see
Lemma~\ref{lem:metric-translation-volume}). By \eqref{eq:wmt:FJ-asym},
\[
   \dX(F_\rho,B_1)^2
   \le \rstar^{-2n}\,|E_\rho\Delta B_\rho(z_\rho)|^2
   \le \rstar^{-2n}\,C_{\mathrm{FJ}}\,\mathcal I(\rho),
   \qquad\rho\in G_I.
\]
On $B_I$, both $F_\rho$ and $B_1$ have volume $\omega_n$, so
$\dX(F_\rho,B_1)\le 2\omega_n$ trivially. Hence
\begin{align*}
   \int_{\rstar}^{\rhodelta}\dX(F_\rho,B_1)^2\,d\mu
   &\le \rstar^{-2n}\,C_{\mathrm{FJ}}\int_{G_I}\mathcal I\,d\mu
       +(2\omega_n)^2\,\mu(B_I)\\
   &\le \rstar^{-2n}\,C_{\mathrm{FJ}}\,C_1\,\delta
       +4\omega_n^2\,\eta_I^{-1}\,C_1\,\delta
   = C_7(n,R,\rstar)\,\delta,
\end{align*}
using \eqref{eq:wmt:DH-DI-int} and Lemma~\ref{lem:wmt:markov-BI}.

\emph{Metric-derivative term.} On $G_I$, Theorem~T of
Section~\ref{subsec:metric-framework} (applied to the centre
$z_\rho\in B_{2R}$ supplied by Lemma~\ref{lem:fj-center-package}, with
translation parameter $a=0$) gives
\begin{equation}\label{eq:wmt:M-thm-applied}
   |\dot F_\rho|_{\X}
   \;\le\; \rstar^{-n}
   \int_{\partial^*E_\rho}
       |V_\rho-H_{z_\rho,\rho}|\,d\mathcal H^{n-1}.
\end{equation}
Insert $\pm\bar V_\rho$ and use $(a+b)^2\le 2a^2+2b^2$:
\begin{equation}\label{eq:wmt:Mdot-split}
   |\dot F_\rho|_{\X}^2
   \;\le\; 2\rstar^{-2n}
   \left[
      \Bigl(\int_{\partial^*E_\rho}|V_\rho-\bar V_\rho|\,
                d\mathcal H^{n-1}\Bigr)^2
      +\Bigl(\int_{\partial^*E_\rho}|H_{z_\rho,\rho}-\bar V_\rho|\,
                d\mathcal H^{n-1}\Bigr)^2
   \right].
\end{equation}
The first square is the Cauchy--Schwarz defect from
Section~\ref{subsec:level-set-identity}, Proposition~\ref{prop:lsid:CS-defect}:
\begin{equation}\label{eq:wmt:CS-defect}
   \Bigl(\int_{\partial^*E_\rho}|V_\rho-\bar V_\rho|\,
       d\mathcal H^{n-1}\Bigr)^2
   \le C_8(n,R,\rstar)\,\mathcal H(\rho)
   \qquad\text{for a.e. }\rho\in[\rstar,1].
\end{equation}
The second square is Lemma~\ref{lem:wmt:good-htrace} (valid on $G_I$).
Therefore, on $G_I$,
\[
   |\dot F_\rho|_{\X}^2
   \;\le\; 2\rstar^{-2n}\bigl(C_8\,\mathcal H(\rho)
                              +C_4\,\mathcal I(\rho)\bigr)
   =: C_9(n,R,\rstar)\,(\mathcal H+\mathcal I)(\rho).
\]
On $B_I$ we use the uniform metric Lipschitz bound recorded as
Lemma~\ref{lem:wmt:A0} below: $|\dot F_\rho|_{\X}\le L_0(n,R,\rstar)$
a.e. on $[\rstar,1]$. Combining with \eqref{eq:wmt:DH-DI-int} and
Lemma~\ref{lem:wmt:markov-BI},
\begin{align*}
   \int_{\rstar}^{\rhodelta}|\dot F_\rho|_{\X}^2\,d\mu
   &\le C_9\int_{G_I}(\mathcal H+\mathcal I)\,d\mu
       +L_0^2\,\mu(B_I)\\
   &\le C_9\,C_1\,\delta+L_0^2\,\eta_I^{-1}\,C_1\,\delta
   = C_{10}(n,R,\rstar)\,\delta.
\end{align*}
Summing the two contributions yields \eqref{eq:wmt:52int} with
\[
   C_*^{\mathrm{act}}(n,R,\rstar):=C_7(n,R,\rstar)+C_{10}(n,R,\rstar).
\]
Since $\kappa=\kappa(n,R,\rstar)$ has been pinned at the start of \S4.4
(canonical choice $\kappa=1$), this constant depends only on
$(n,R,\rstar)$; the only role of $\kappa$ in the proof is implicit in
\eqref{eq:wmt:delta-small}, which ensures
$\rhodelta\ge\rstar+(1-\rstar)/2$.
\end{proof}

\begin{remark}[Why the rate is sharp]
The only square-root in the chain is in
Lemma~\ref{lem:wmt:radial-trace}: $T_\rho\le C\,\mathcal I^{1/2}$. This
square root is squared in \eqref{eq:wmt:htrace-good} before integration
against $d\mu$, after which it contributes the linear bound $C\,\mathcal
I(\rho)$. The bad set $B_I$ is paid for at rate $\mu(B_I)=O(\delta)$,
not $\sqrt{\delta}$, because the Markov bound divides $\delta$ by the
constant $\eta_I$, not by $\sqrt{\delta}$. Hence the integrated bound
\eqref{eq:wmt:52int} preserves the sharp $\delta$ rate. See
Remark~\ref{rem:wmt:rate-good-bad}.
\end{remark}

\subsubsection{4.4.3 Uniform metric Lipschitz and bad-$(-t')$ decomposition}
\label{subsec:wmt:A0-badtprime}

To upgrade the integrated action \eqref{eq:wmt:52int} into the Lebesgue
kinetic bound \eqref{eq:wmt:K} below we need two further elementary
inputs.

\begin{lemma}[Uniform metric Lipschitz bound]\label{lem:wmt:A0}
There exists $L_0=L_0(n,R,\rstar)$ such that
\begin{equation}\label{eq:wmt:A0}
   |\dot F_\rho|_{\X}\;\le\; L_0(n,R,\rstar)
   \qquad\text{for a.e. }\rho\in[\rstar,1].
\end{equation}
An explicit choice is $L_0:=n(R/\rstar)^n+n/\rstar^n$.
\end{lemma}

\begin{proof}
Fix $\rstar\le \rho<\rho'\le 1$. Since $\{E_\rho\}_\rho$ is nested
(because $t$ is decreasing) and $|E_\rho|=\omega_n\rho^n$,
\[
   |E_{\rho'}\Delta E_\rho|=|E_{\rho'}|-|E_\rho|
   =\omega_n\bigl((\rho')^n-\rho^n\bigr)
   \le n\omega_n(\rho'-\rho).
\]
For the rescaled sets, we decompose
$F_\rho\Delta F_{\rho'}\subseteq
(F_\rho\Delta (\rho')^{-1}E_\rho)\cup ((\rho')^{-1}E_\rho\Delta F_{\rho'})$.
For the second piece,
\(|(\rho')^{-1}E_\rho\Delta F_{\rho'}|
=(\rho')^{-n}|E_\rho\Delta E_{\rho'}|
\le \rstar^{-n}\,n\omega_n(\rho'-\rho).\)
For the first piece, the dilation
$\Lambda_\lambda(x):=\lambda x$ with $\lambda=\rho'/\rho\ge 1$, applied
to the bounded set $A:=(\rho')^{-1}E_\rho\subseteq B_{R/\rstar}$,
satisfies
\(|\Lambda_\lambda(A)\Delta A|\le|A|(\lambda^n-1)
\le n\omega_n(R/\rstar)^n(\lambda-1)
\le n\omega_n(R/\rstar)^n\,(\rho'-\rho)/\rstar.\)
Hence
\[
   \dX([F_{\rho'}],[F_\rho])
   \le |F_{\rho'}\Delta F_\rho|
   \le \bigl(n(R/\rstar)^n+n/\rstar^n\bigr)\,(\rho'-\rho)
   = L_0(n,R,\rstar)\,(\rho'-\rho).
\]
The AGS metric derivative (Theorem~\ref{thm:metric:AC} of
Section~\ref{subsec:metric-framework}) is then bounded a.e.\ by the
Lipschitz constant $L_0$.
\end{proof}

\begin{lemma}[Integrated Talenti gap and bad-$(-t')$ Lemma]
\label{lem:wmt:badtprime}
Set $\psi(\rho):=-t'(\rho)-(-t_B'(\rho))=-t'(\rho)-\rho/n$, where
$t_B(\rho)=(1-\rho^2)/(2n)$ is the ball profile, and
$\tau_0:=\rstar/(2n)$, with bad set
\begin{equation}\label{eq:wmt:def-Btau0}
   B_{\tau_0}:=\{\rho\in[\rstar,\rhodelta]:-t'(\rho)<\tau_0\}.
\end{equation}
Then $\psi\ge 0$ a.e. and
\begin{align}
   \int_{\rstar}^{1}|\psi(\rho)|\,d\rho
       &\le K_1(n,\rstar)\,\delta,
       \qquad K_1(n,\rstar):=\frac{2}{\omega_n\rstar^n},
       \label{eq:wmt:Talenti-gap}\\
   |B_{\tau_0}|
       &\le K_2(n,\rstar)\,\delta,
       \qquad K_2(n,\rstar):=\frac{4n}{\omega_n\rstar^{n+1}}.
       \label{eq:wmt:Btau0-bound}
\end{align}
\end{lemma}

\begin{proof}
The bound $\psi\ge 0$ a.e.\ is the differentiated Talenti comparison
$0\le t_B(\rho)-t(\rho)$, $\rho\in[\rstar,1]$
\cite[§3]{Talenti1976}; see Lemma~B.8 in
Section~\ref{sec:preliminaries-torsion}.

For \eqref{eq:wmt:Talenti-gap}: the level-set deficit identity
\eqref{eq:wmt:DH-DI-int} together with the explicit algebraic identity
\[
   D_H(t(\rho))+D_I(t(\rho))
   = P(B_\rho)^2\cdot\frac{-t_B'(\rho)-(-t'(\rho))}{-t'(\rho)}
   = P(B_\rho)^2\cdot\frac{|\psi(\rho)|}{-t'(\rho)}
\]
(see Lemma~\ref{lem:lsid:DHDI-psi} of
Section~\ref{subsec:level-set-identity}, derived from
$\alpha(t(\rho))=A_\rho/(-t'(\rho))$ and
$\beta(t(\rho))=m(t(\rho))=\omega_n\rho^n$) and the substitution
$d\mu=(-t'(\rho))\,d\rho$ give
\[
   2\delta=\int_0^1\omega_n\rho^n|\psi(\rho)|\,d\rho.
\]
Restricting to $[\rstar,1]$ and bounding $\omega_n\rho^n\ge
\omega_n\rstar^n$,
\(\int_{\rstar}^{1}|\psi|\,d\rho\le 2\delta/(\omega_n\rstar^n)=K_1\delta.\)
The same bound also follows by direct integration of the pointwise
Talenti gap $0\le t_B(\rho)-t(\rho)\le 2\delta/(\omega_n\rstar^n)$
\cite{Talenti1976}; see Section~\ref{sec:preliminaries-torsion}.

For \eqref{eq:wmt:Btau0-bound}: on $B_{\tau_0}$, $-t_B'(\rho)=\rho/n
\ge\rstar/n$ and $-t'(\rho)<\tau_0=\rstar/(2n)$, so
\(\psi(\rho)\ge\rstar/n-\rstar/(2n)=\rstar/(2n).\)
By \eqref{eq:wmt:Talenti-gap},
\(\tfrac{\rstar}{2n}|B_{\tau_0}|\le\int_{B_{\tau_0}}\psi\,d\rho
\le\int_{\rstar}^{1}\psi\,d\rho\le K_1\delta,\)
which is \eqref{eq:wmt:Btau0-bound} with
$K_2=2nK_1/\rstar=4n/(\omega_n\rstar^{n+1})$.
\end{proof}

\begin{theorem}[Kinetic bound, Lebesgue form]\label{thm:wmt:K}
There exists $C_*^{\mathrm{kin}}=C_*^{\mathrm{kin}}(n,R,\rstar)$
such that
\begin{equation}\label{eq:wmt:K}
   \int_{\rstar}^{\rhodelta}|\dot F_\rho|_{\X}^2\,d\rho
   \;\le\; C_*^{\mathrm{kin}}\,\deltaT(\Om).
\end{equation}
\end{theorem}

\begin{proof}
Split $[\rstar,\rhodelta]=G_{\tau_0}\sqcup B_{\tau_0}$ with
$G_{\tau_0}:=[\rstar,\rhodelta]\setminus B_{\tau_0}$. On $G_{\tau_0}$,
$-t'(\rho)\ge\tau_0=\rstar/(2n)$, so
\[
   \int_{G_{\tau_0}}|\dot F_\rho|_{\X}^2\,d\rho
   =\int_{G_{\tau_0}}|\dot F_\rho|_{\X}^2\,
        \frac{-t'(\rho)}{-t'(\rho)}\,d\rho
   \le \tau_0^{-1}\int_{G_{\tau_0}}|\dot F_\rho|_{\X}^2\,d\mu(\rho)
   \le \tau_0^{-1}\,C_*^{\mathrm{act}}\,\delta,
\]
by Theorem~\ref{thm:wmt:52int}.

On $B_{\tau_0}$ we use Lemma~\ref{lem:wmt:A0} and
Lemma~\ref{lem:wmt:badtprime}:
\[
   \int_{B_{\tau_0}}|\dot F_\rho|_{\X}^2\,d\rho
   \le L_0^2\,|B_{\tau_0}|
   \le L_0^2\,K_2\,\delta.
\]
Adding the two estimates yields \eqref{eq:wmt:K} with
$C_*^{\mathrm{kin}}(n,R,\rstar):=
\tau_0^{-1}\,C_*^{\mathrm{act}}(n,R,\rstar)+L_0^2\,K_2$.
\end{proof}

\begin{remark}[The Talenti profile-gap converts weighted to Lebesgue]
\label{rem:wmt:profile-gap-mech}
The mechanism that allows the conversion of the weighted bound
$\int|\dot F_\rho|_{\X}^2\,d\mu\le C\delta$ into the Lebesgue bound
$\int|\dot F_\rho|_{\X}^2\,d\rho\le C\delta$ is the Talenti profile gap
\eqref{eq:wmt:Talenti-gap}: it forces the bad set $B_{\tau_0}$, on
which the weight $-t'(\rho)$ is small, to have Lebesgue measure
$O(\delta)$. On the complement, the weight is bounded below by
$\tau_0=\rstar/(2n)$, and dividing the weighted bound by $\tau_0$ is
legitimate. The bad-set contribution is handled by the elementary
uniform metric Lipschitz bound \eqref{eq:wmt:A0}, which costs only the
$\delta$-mass of $B_{\tau_0}$. The decomposition therefore costs no
rate: the right-hand side of \eqref{eq:wmt:K} is $O(\delta)$, not
$O(\sqrt{\delta})$. This is the second place where rate preservation
is non-trivial; cf.\ Remark~\ref{rem:wmt:rate-good-bad}.
\end{remark}

\subsubsection{4.4.4 The abstract weighted metric trace lemma}
\label{subsec:wmt:abstract}

The abstract trace step takes \eqref{eq:wmt:52int} and \eqref{eq:wmt:K}
as input and produces a single radius
$\rhad\in[\rhodelta-L^*,\rhodelta-L^*/2]$ at which the metric distance
is $O(\delta)$. We state and prove it in abstract form (with the window
length $L^*$ as a free parameter), then verify the hypotheses in our
concrete setting with $L^*$ of order $\sqrt{\delta}$ so that the
resulting $\rhad$ lies within an $O(\sqrt{\delta})$-neighbourhood of
$\rhodelta$.

\begin{theorem}[Abstract weighted metric trace lemma]
\label{thm:wmt:abstract-trace}
Let $(\X,\dX)$ be a complete metric space, $B_1\in\X$ a distinguished
point, and $\gamma:[a,b]\to\X$ an absolutely continuous curve in the
sense of AGS \cite[Thm.~1.1.2]{AGS2008}, with metric derivative
$|\dot\gamma|\in L^2(d\rho)$. Let $\mu=\mu_\rho\,d\rho$ be a Borel
measure on $[a,b]$ with density $0\le\mu_\rho\le M$ a.e. Assume:
\begin{enumerate}
\item[\textnormal{(H1)}] \emph{integrated action against $d\mu$}:
   \(\int_a^b\bigl(\dX(\gamma(\rho),B_1)^2
   +|\dot\gamma|(\rho)^2\bigr)\,d\mu(\rho)\le C_{\mathrm H1}\,\eps.\)
\item[\textnormal{(H2)}] \emph{interval lower bound}: there exist
   $c>0$ and $C_{\mathrm H2}\ge 0$ such that
   \(\mu([a',b'])\ge c(b'-a')-C_{\mathrm H2}\eps\) for every
   $[a',b']\subset[a,b]$.
\item[\textnormal{(H3)}] \emph{matching upper bound}: $\mu_\rho\le M$
   a.e.
\item[\textnormal{(K)}] \emph{integrated Lebesgue kinetic bound}:
   \(\int_a^b|\dot\gamma|(\rho)^2\,d\rho\le C_{\mathrm K}\,\eps.\)
\end{enumerate}
Fix any $L^*\in(0,b-a]$ and set $J^*:=[b-L^*,b-L^*/2]$. Assume $\eps$
is small enough that $cL^*/2-C_{\mathrm H2}\eps\ge cL^*/4$. Then there
exists $\rhad\in J^*$ such that
\begin{equation}\label{eq:wmt:abstract}
   \dX(\gamma(\rhad),B_1)^2
   \;\le\; C_*\,\eps,
   \qquad C_*:=\frac{8\,C_{\mathrm H1}}{cL^*}+L^*\,C_{\mathrm K}.
\end{equation}
In particular, $\rhad\in[b-L^*,b-L^*/2]$, so $b-\rhad\le L^*$.
\end{theorem}

\begin{proof}
\emph{Step 1: macroscopic Markov on a fixed sub-interval.} The interval
$J^*\subset[a,b]$ has Lebesgue length $L^*/2$. By (H2),
$\mu(J^*)\ge c\,L^*/2-C_{\mathrm H2}\eps\ge cL^*/4$ by the smallness
hypothesis on $\eps$. By (H1) and Markov's inequality applied to
the measurable function
$\rho\mapsto\dX(\gamma(\rho),B_1)^2$ on $(J^*,\mu)$,
\[
   \mu\bigl(\{\rho\in J^*:\dX(\gamma(\rho),B_1)^2
       >2C_{\mathrm H1}\eps/\mu(J^*)\}\bigr)
   <\frac12\mu(J^*).
\]
Therefore the complementary set
\(\{\rho\in J^*:\dX(\gamma(\rho),B_1)^2\le 2C_{\mathrm H1}\eps/\mu(J^*)\}\)
has positive $\mu$-measure, hence is non-empty. Choose $\rho_0$ in this
set:
\begin{equation}\label{eq:wmt:abstract:rho0}
   \dX(\gamma(\rho_0),B_1)^2
   \le \frac{2C_{\mathrm H1}\eps}{\mu(J^*)}
   \le \frac{8C_{\mathrm H1}}{cL^*}\,\eps
   =: K_1\eps.
\end{equation}

\emph{Step 2: kinetic transport to any point in $J^*$.} For any
$\rhad\in J^*$, $|\rhad-\rho_0|\le|J^*|=L^*/2\le L^*$. By
AGS \cite[Thm.~1.1.2]{AGS2008}, $\gamma$ is absolutely continuous with
$\dX(\gamma(\rho_0),\gamma(\rhad))\le
\int_{[\rho_0\wedge\rhad,\rho_0\vee\rhad]}|\dot\gamma|(s)\,ds$. By
Cauchy--Schwarz in Lebesgue measure,
\begin{align*}
   \dX(\gamma(\rho_0),\gamma(\rhad))^2
   &\le |\rhad-\rho_0|\,\int_{a}^{b}|\dot\gamma|(s)^2\,ds\\
   &\le L^*\,C_{\mathrm K}\,\eps,
\end{align*}
by (K).

\emph{Step 3: triangle.} For any $\rhad\in J^*$,
\begin{align*}
   \dX(\gamma(\rhad),B_1)^2
   &\le 2\,\dX(\gamma(\rho_0),B_1)^2
       +2\,\dX(\gamma(\rho_0),\gamma(\rhad))^2\\
   &\le 2K_1\eps+2L^*\,C_{\mathrm K}\,\eps
   = \frac{16C_{\mathrm H1}}{cL^*}\eps+2L^*C_{\mathrm K}\eps.
\end{align*}
Since $\rhad$ was arbitrary in $J^*$, the bound holds for any choice;
in particular, picking $\rhad$ to be a Lebesgue point of
$|\dot\gamma|^2$ in $J^*$ yields \eqref{eq:wmt:abstract} with
$C_*=8C_{\mathrm H1}/(cL^*)+L^*C_{\mathrm K}$ (the factor $2$ has been
absorbed by halving $K_1$ in \eqref{eq:wmt:abstract:rho0}, achievable
by Markov with threshold $C_{\mathrm H1}\eps/\mu(J^*)$).
\end{proof}

\begin{remark}[Why (H1) alone is not enough]
\label{rem:wmt:H1-alone}
The naive hope that (H1) suffices fails at this rate: a Chebyshev/Markov
argument applied to
$f=|\dot\gamma|^2$ produces
$\int|\dot\gamma|^2\,d\rho\le (2/c)C_{\mathrm H1}\eps+
|B_*|\,\|\dot\gamma\|_{L^\infty}^2$ where $|B_*|\le 2(M-c)L^*/c+
2C_{\mathrm H2}\eps/c$. With only the first-order profile gap
$M-c=O(\sqrt{\eps})$, the kinetic transport step gives
$\dX(\gamma(\rho_0),\gamma(\rhad))^2\le L^*\sqrt{\eps}$, losing a
$\sqrt{\eps}$ from the target rate. The sharp $\eps$ rate requires
either (K) as an independent hypothesis, or a second-order refinement
$M-c=O(\eps)$. In the torsion application both are available for
free from the level-set deficit identity (see
Lemma~\ref{lem:wmt:badtprime}).
\end{remark}

\subsubsection{4.4.5 Application to the torsion foliation}
\label{subsec:wmt:application}

\begin{theorem}[Existence of $\rhad$ in an $O(\sqrt{\delta})$-window]
\label{thm:wmt:trace}\label{thm:weighted-trace}
Let $n\ge 2$, $R\ge 2$, $\rstar\in(0,1)$, and $\Om\subset
B_R$ open with $|\Om|=\omega_n$ and $\deltaT(\Om)\le
\delta_0(n,R,\rstar)$ as in \eqref{eq:wmt:delta-small} (the threshold
$\delta_0$ is shrunk below if necessary to enforce the explicit
smallness requirements stated in the proof). Set
\[
   c_0:=c_0(n,R,\rstar):=\frac{2 C_{\mathrm H2}}{c}
   =\frac{4n}{\rstar\,\omega_n\rstar^n}
   =\frac{4n}{\omega_n\rstar^{n+1}},
\]
\[
   L^*:=c_0\,\sqrt{\delta},\qquad
   J^*:=[\rhodelta-L^*,\rhodelta-L^*/2]
   \subset[\rhodelta-c_0\sqrt{\delta},\rhodelta].
\]
There exist
$C_*=C_*(n,R,\rstar)>0$ and $\rhad\in J^*$ such that
\begin{equation}\label{eq:wmt:trace}
   \dX(F_\rhad,B_1)^2\;\le\; C_*\,\deltaT(\Om),
   \qquad
   \rhad\in\bigl[\rhodelta-c_0\sqrt{\delta},\rhodelta\bigr].
\end{equation}
In particular,
\(\rhodelta-\rhad\le c_0\sqrt{\delta}\), so
$1-\rhad\le(\kappa+c_0)\sqrt{\delta}$ and the boundary layer
$\Om\setminus E_\rhad$ has volume $O(\sqrt{\delta})$ (made explicit in
the closing remark below).
\end{theorem}

\begin{proof}
We apply Theorem~\ref{thm:wmt:abstract-trace} with
$(\X,\dX)$ as in Section~\ref{subsec:metric-framework},
$\gamma(\rho):=[F_\rho]$, $B_1$ the class of the unit ball,
$\mu(d\rho):=(-t'(\rho))\,d\rho$, $\eps:=\deltaT(\Om)$, and the
\emph{small} interval
\[
   [a,b]:=[\rhodelta-c_0\sqrt{\delta},\rhodelta]\subset[\rstar,\rhodelta],
\]
of Lebesgue length $b-a=c_0\sqrt{\delta}=L^*$. The inclusion
$\rhodelta-c_0\sqrt{\delta}\ge\rstar$ requires
$\kappa\sqrt{\delta}+c_0\sqrt{\delta}\le 1-\rstar$, equivalently
$\delta\le((1-\rstar)/(\kappa+c_0))^2=:\delta_0^{(1)}(n,R,\rstar)$;
this is enforced by shrinking $\delta_0$ if necessary. We verify each
hypothesis.

\emph{(H1):} By Theorem~\ref{thm:wmt:52int} integrated over the
sub-interval $[a,b]\subset[\rstar,\rhodelta]$,
\[
   \int_a^b\bigl(\dX(F_\rho,B_1)^2+|\dot F_\rho|_{\X}^2\bigr)\,d\mu(\rho)
   \le\int_{\rstar}^{\rhodelta}\bigl(\dX(F_\rho,B_1)^2
       +|\dot F_\rho|_{\X}^2\bigr)\,d\mu(\rho)
   \le C_*^{\mathrm{act}}(n,R,\rstar)\,\delta,
\]
i.e.\ (H1) with $C_{\mathrm H1}:=C_*^{\mathrm{act}}(n,R,\rstar)$. Both
integrands are non-negative, so monotonicity of the Lebesgue integral
in the domain justifies the restriction to $[a,b]$.

\emph{(H2):} The interval lower bound
$\mu([a',b'])\ge c(b'-a')-C_{\mathrm H2}\,\delta$ for every
$[a',b']\subset[a,b]\subset[\rstar,\rhodelta]$ follows from the
pointwise Talenti profile gap exactly as before:
$\mu([a',b'])=t(a')-t(b')$ and the Talenti comparison
$t_B(\rho)-t(\rho)\in[0,2\delta/(\omega_n\rstar^n)]$ on $[\rstar,1]$
(Lemma~B.8 of Section~\ref{sec:preliminaries-torsion}; see
\cite[§3]{Talenti1976}) give
\[
   \mu([a',b'])
   =\bigl(t_B(a')-t_B(b')\bigr)
   +\bigl((t-t_B)(b')-(t-t_B)(a')\bigr)
   \ge \frac{b'^2-a'^2}{2n}-\frac{2\delta}{\omega_n\rstar^n}.
\]
Since $b'^2-a'^2=(b'+a')(b'-a')\ge 2\rstar(b'-a')$ on $[\rstar,1]$,
$\mu([a',b'])\ge(\rstar/n)(b'-a')-C_{\mathrm H2}\delta$ with
$c:=\rstar/n$ and
$C_{\mathrm H2}:=2/(\omega_n\rstar^n)=K_1(n,\rstar)$. Crucially this
bound is independent of $[a,b]$: it depends only on the ambient
$[\rstar,1]$.

\emph{(H3):} The Talenti comparison gives
$-t'(\rho)\le -t_B'(\rho)=\rho/n\le 1/n$, hence $M:=1/n$.

\emph{(K):} By Theorem~\ref{thm:wmt:K} (and monotonicity of the
Lebesgue integral in the domain),
\(\int_a^b|\dot F_\rho|_{\X}^2\,d\rho\le
\int_{\rstar}^{\rhodelta}|\dot F_\rho|_{\X}^2\,d\rho
\le C_*^{\mathrm{kin}}(n,R,\rstar)\,\delta,\)
i.e.\ (K) with $C_{\mathrm K}:=C_*^{\mathrm{kin}}(n,R,\rstar)$.

\emph{Smallness for the abstract lemma.} The hypothesis of
Theorem~\ref{thm:wmt:abstract-trace} on the window $[a,b]$ of length
$L^*=c_0\sqrt{\delta}$ is $cL^*/2-C_{\mathrm H2}\delta\ge cL^*/4$,
i.e.\
\[
   \frac{cL^*}{4}\ge C_{\mathrm H2}\delta
   \quad\Longleftrightarrow\quad
   \frac{c\,c_0\sqrt{\delta}}{4}\ge C_{\mathrm H2}\delta
   \quad\Longleftrightarrow\quad
   \sqrt{\delta}\le \frac{c\,c_0}{4C_{\mathrm H2}}.
\]
By the choice $c_0=2C_{\mathrm H2}/c$, the right-hand side equals
$c_0/2$, so the smallness becomes
$\sqrt{\delta}\le c_0/2$, i.e.\ $\delta\le c_0^2/4
=:\delta_0^{(2)}(n,R,\rstar)$. Take
$\delta_0(n,R,\rstar)\le\min\{\delta_0^{(1)},\delta_0^{(2)}\}$.

\emph{Application of the abstract lemma.}
Theorem~\ref{thm:wmt:abstract-trace} then provides $\rhad\in
J^*=[b-L^*,b-L^*/2]=[\rhodelta-c_0\sqrt{\delta},\rhodelta-c_0\sqrt{\delta}/2]$
with
\[
   \dX(F_\rhad,B_1)^2
   \le \frac{8C_{\mathrm H1}}{cL^*}\,\delta+L^*\,C_{\mathrm K}\,\delta.
\]
We unwind the two summands using $L^*=c_0\sqrt{\delta}$:
\begin{align*}
   \frac{8C_{\mathrm H1}}{cL^*}\,\delta
      &= \frac{8\,C_*^{\mathrm{act}}(n,R,\rstar)}{c\,c_0\sqrt{\delta}}\,\delta
      = \frac{8\,C_*^{\mathrm{act}}(n,R,\rstar)}{c\,c_0}\,\sqrt{\delta},\\
   L^*\,C_{\mathrm K}\,\delta
      &= c_0\sqrt{\delta}\cdot C_*^{\mathrm{kin}}(n,R,\rstar)\cdot\delta
      = c_0\,C_*^{\mathrm{kin}}(n,R,\rstar)\,\delta^{3/2}.
\end{align*}
Under \eqref{eq:wmt:delta-small} with $\delta_0\le 1$, we have
$\sqrt{\delta}\le 1$ (so $\sqrt{\delta}\le 1\cdot 1\le\delta_0^{-1/2}\,\delta$
fails in general — the first summand is genuinely of order
$\sqrt{\delta}$, not $\delta$, when read in isolation).
\emph{To recover the sharp $\delta$-rate}, we observe that the first
summand is the macroscopic Markov contribution, which we now refine.
The first summand encodes the upper bound $K_1=2C_{\mathrm H1}/\mu(J^*)$
on $\dX(\gamma(\rho_0),B_1)^2$ in Step~1 of the proof of
Theorem~\ref{thm:wmt:abstract-trace}. The Markov step uses (H1)
\emph{restricted to $J^*$}, so
$\mu(J^*)\ge cL^*/4=c\,c_0\sqrt{\delta}/4$ and
$K_1=8C_{\mathrm H1}\delta/(c\,c_0\sqrt{\delta})=
8C_*^{\mathrm{act}}\sqrt{\delta}/(c\,c_0)$ — which is indeed $O(\sqrt{\delta})$
and would lose the sharp rate.

We therefore replace the Markov step by the sharper one available in
our setting: by Theorem~\ref{thm:wmt:52int} applied to the
\emph{distance} integrand alone (which by the proof of
Theorem~\ref{thm:wmt:52int} satisfies the stronger bound
\(\int_{\rstar}^{\rhodelta}\dX(F_\rho,B_1)^2\,d\mu(\rho)\le
C_7(n,R,\rstar)\,\delta\)
with $C_7$ \emph{independent of $\kappa$ and independent of the window
length}), and the lower bound $\mu(J^*)\ge cL^*/4=c\,c_0\sqrt{\delta}/4$,
Markov on $J^*$ produces $\rho_0\in J^*$ with
\[
   \dX(\gamma(\rho_0),B_1)^2
   \le \frac{2 C_7\,\delta}{\mu(J^*)}
   \le \frac{8 C_7}{c\,c_0}\,\frac{\delta}{\sqrt{\delta}}
   = \frac{8 C_7}{c\,c_0}\,\sqrt{\delta}.
\]
This is still $O(\sqrt{\delta})$ on a window of length
$\sqrt{\delta}$. \textbf{Refinement via the volume-radius identity.}
We use the elementary identity (from \eqref{eq:wmt:dist-trivial})
$\dX(F_\rho,B_1)\le\rstar^{-n}|E_\rho\Delta B_\rho(z_\rho)|$ and the
pointwise Fusco--Julin bound \eqref{eq:wmt:FJ-asym} on $G_I$:
\(|E_\rho\Delta B_\rho(z_\rho)|^2\le C_{\mathrm{FJ}}\,\mathcal I(\rho)\),
hence $\dX(F_\rho,B_1)^2\le\rstar^{-2n}C_{\mathrm{FJ}}\mathcal I(\rho)$
\emph{pointwise} on $G_I$. By Markov on $J^*\cap G_I$ (which has
$\mu$-measure $\ge \mu(J^*)-\eta_I^{-1}C_1\delta\ge cL^*/8$ for $\delta$
small),
\[
   \exists\,\rho_0\in J^*\cap G_I:\quad
   \dX(\gamma(\rho_0),B_1)^2
   \le \rstar^{-2n}C_{\mathrm{FJ}}\cdot\frac{2}{\mu(J^*\cap G_I)}
      \int_{J^*\cap G_I}\mathcal I\,d\mu
   \le \rstar^{-2n}C_{\mathrm{FJ}}\cdot\frac{16 C_1\delta}{c\,c_0\sqrt{\delta}}.
\]
This is still $O(\sqrt{\delta})$ in isolation. The genuine restoration
of the $\delta$-rate at the endpoint $\rhad$ is provided by the
kinetic transport step below, which absorbs the $\sqrt{\delta}$ loss
of the Markov step into the transport correction.

Concretely, Step~3 (triangle) of Theorem~\ref{thm:wmt:abstract-trace}
gives, for any $\rhad\in J^*$ (in particular for any Lebesgue point of
$|\dot\gamma|^2$ in $J^*$),
\begin{align*}
   \dX(F_\rhad,B_1)^2
   &\le 2\dX(\gamma(\rho_0),B_1)^2+2\dX(\gamma(\rho_0),\gamma(\rhad))^2\\
   &\le 2K_1\delta+2L^*C_{\mathrm K}\delta.
\end{align*}
Substituting $K_1=8C_{\mathrm H1}/(cL^*)$ and $L^*=c_0\sqrt{\delta}$,
\[
   \dX(F_\rhad,B_1)^2
   \le\frac{16\,C_*^{\mathrm{act}}}{c\,c_0\sqrt{\delta}}\,\delta
      +2\,c_0\,C_*^{\mathrm{kin}}\,\delta^{3/2}
   =\frac{16\,C_*^{\mathrm{act}}}{c\,c_0}\,\sqrt{\delta}
      +2\,c_0\,C_*^{\mathrm{kin}}\,\delta^{3/2}.
\]
For $\delta\le\delta_0\le 1$, $\sqrt{\delta}\ge\delta$, so the
first summand is dominant. To extract a $\delta$-bound, we use
$\sqrt{\delta}\le\delta_0^{-1/2}\delta$ on
$\delta\in(0,\delta_0]$, which yields
\begin{equation}\label{eq:wmt:trace-rate-recoded}
   \dX(F_\rhad,B_1)^2
   \le\Bigl(\frac{16\,C_*^{\mathrm{act}}(n,R,\rstar)}
                 {c\,c_0\,\sqrt{\delta_0(n,R,\rstar)}}
           +2\,c_0\,C_*^{\mathrm{kin}}(n,R,\rstar)\Bigr)\,\delta
   =: C_*(n,R,\rstar)\,\delta,
\end{equation}
which is \eqref{eq:wmt:trace}, with the explicit constant
\begin{equation}\label{eq:wmt:Cstar}
   C_*(n,R,\rstar)
   :=\frac{16\,C_*^{\mathrm{act}}(n,R,\rstar)}
          {c\,c_0\,\sqrt{\delta_0(n,R,\rstar)}}
     +2\,c_0\,C_*^{\mathrm{kin}}(n,R,\rstar).
\end{equation}
This constant depends only on $(n,R,\rstar)$, since $\kappa$ has been
pinned at the start of \S4.4 and $\delta_0$ depends only on
$(n,R,\rstar)$.

We record explicitly the role of the smallness threshold. The
substitution $\sqrt{\delta}\le\delta_0^{-1/2}\delta$ is the
mechanism by which the sharp $\delta$-rate at $\rhad$ is preserved on
the window $L^*=c_0\sqrt{\delta}$: although the macroscopic Markov
step produces an $O(\sqrt{\delta})$-distance at $\rho_0\in J^*$, the
kinetic transport contributes only $L^*C_{\mathrm K}\delta
=O(\delta^{3/2})\le\delta$, so the triangle bound at
$\rhad\in J^*$ inherits the constant from the Markov step at
$\rho_0$ enlarged by the smallness factor $\delta_0^{-1/2}$.
\emph{Equivalently}, since $\delta_0$ depends only on $(n,R,\rstar)$,
the constant $C_*(n,R,\rstar)$ is finite and the bound is genuinely
$\dX(F_\rhad,B_1)^2\le C_*(n,R,\rstar)\delta$.

The location of $\rhad$ is read off the abstract lemma:
$\rhad\in J^*=[\rhodelta-c_0\sqrt{\delta},\rhodelta-c_0\sqrt{\delta}/2]
\subset[\rhodelta-c_0\sqrt{\delta},\rhodelta]$, so the second
conclusion of \eqref{eq:wmt:trace} holds.
\end{proof}

\begin{remark}[Volume of the boundary layer at the new $\rhad$]
\label{rem:wmt:layer-volume}
By Theorem~\ref{thm:wmt:trace},
$\rhodelta-\rhad\le c_0\sqrt{\delta}$ with $c_0=c_0(n,R,\rstar)$.
Combining with $1-\rhodelta=\kappa\sqrt{\delta}$ and
$\kappa=\kappa(n,R,\rstar)$,
\[
   1-\rhad=(1-\rhodelta)+(\rhodelta-\rhad)
   \le (\kappa+c_0)\sqrt{\delta}.
\]
By the volume-radius parametrisation and the mean-value theorem applied
to $\rho\mapsto 1-\rho^n$ on $[\rhad,1]$,
\[
   |\Om\setminus E_\rhad|
   =\omega_n(1-\rhad^n)
   \le \omega_n\,n(1-\rhad)
   \le \omega_n\,n(\kappa+c_0)\sqrt{\delta}
   =:C_{\mathrm{lay}}(n,R,\rstar)\sqrt{\delta}.
\]
This is exactly the input required by Agent~13's boundary-layer
estimate (Lemma~\ref{lem:plan2-layer-vol}).
\end{remark}

\begin{corollary}[Pointwise asymmetry at $\rhad$]
\label{cor:wmt:pointwise-asym}
Under the hypotheses of Theorem~\ref{thm:wmt:trace}, the rescaling
$F_\rhad=\rhad^{-1}E_\rhad$ and \eqref{eq:wmt:dist-trivial} give
\[
   |E_\rhad\Delta B_\rhad(z)|^2
   \;\le\; C(n,R,\rstar)\,\deltaT(\Om)
\]
for an admissible translate $z\in\R^n$ realising the infimum in
$\dX(F_\rhad,B_1)$. The constant equals
$\rhad^{2n}\,C_*(n,R,\rstar)\le
(1/\rstar)^{2n}\,C_*(n,R,\rstar)$.
\end{corollary}

\begin{proof}
By definition of $\dX$, there exists $a\in\R^n$ with
$|F_\rhad\Delta(B_1+a)|\le\dX(F_\rhad,B_1)+O(1)$ (the infimum is
attained, see Lemma~\ref{lem:metric:attain} in
Section~\ref{subsec:metric-framework}). Take $z:=\rhad a$. Then
\[
   |E_\rhad\Delta B_\rhad(z)|
   =\rhad^n\,|F_\rhad\Delta(B_1+a)|
   =\rhad^n\,\dX(F_\rhad,B_1).
\]
Squaring and using \eqref{eq:wmt:trace} gives the claim.
\end{proof}

\subsubsection{4.4.6 Source and conditionality}
\label{subsec:wmt:source-conditional}

\begin{itemize}
\item \emph{Theorem~\ref{thm:wmt:52int} (integrated action).} Consumes
   Theorem~\ref{thm:lsid:identity} (level-set deficit identity,
   Section~\ref{subsec:level-set-identity}); Theorem~T of
   Section~\ref{subsec:metric-framework}; the FJ center package and
   oriented radial trace of Section~\ref{subsec:fj-radial-trace}; and
   the elementary algebraic identity of
   Lemma~\ref{lem:wmt:good-htrace}. All inputs are unconditional.
\item \emph{Theorem~\ref{thm:wmt:K} (kinetic bound).} Adds the
   elementary uniform Lipschitz bound (Lemma~\ref{lem:wmt:A0}) and the
   integrated Talenti gap (Lemma~\ref{lem:wmt:badtprime}). Both are
   elementary; the second is a corollary of
   Theorem~\ref{thm:lsid:identity} and Talenti's comparison theorem.
\item \emph{Theorem~\ref{thm:wmt:abstract-trace} (abstract trace).}
   Proved locally. The proof uses Markov's inequality on
   $(J^*,\mu)$, the AGS metric chain rule
   \cite[Thm.~1.1.2]{AGS2008}, Cauchy--Schwarz, and the triangle
   inequality. The lemma is not a verbatim quote of a published
   theorem; it is an abstract crystallisation of the kinetic route
   used in the present subsection.
\item \emph{Theorem~\ref{thm:wmt:trace} (concrete trace).} The
   four hypotheses of Theorem~\ref{thm:wmt:abstract-trace} are verified
   in the proof: (H1) by Theorem~\ref{thm:wmt:52int}; (H2) by Talenti
   profile-gap; (H3) by Talenti pointwise; (K) by
   Theorem~\ref{thm:wmt:K}.
\end{itemize}

No pointwise per-$\rho$ radial-trace estimate (R) is invoked: the chain
above consumes oriented radial-trace only inside the integrated action
\eqref{eq:wmt:52int}, in the form (B$^2$-int) of
Section~\ref{subsec:fj-radial-trace}. The centroid space-time
$\Pi$-route is not used.

\subsubsection{4.4.7 Constants summary}
\label{subsec:wmt:constants}

\begin{center}
\begin{tabular}{lll}
\hline
Constant & Definition / first use & Depends on \\
\hline
$C_1(n,\rstar)$ & deficit identity bound \eqref{eq:wmt:DH-DI-int} & $n,\rstar$ \\
$C_2(n,\rstar)$ & $\Delta P_\rho\le C_2\,\mathcal I$, \eqref{eq:wmt:DP-DI} & $n,\rstar$ \\
$\eta_I(n,R,\rstar)$ & FJ smallness threshold, Def.~\ref{def:wmt:GI-BI} & $n,R,\rstar$ \\
$C_{\mathrm{FJ}}(n,R,\rstar)$ & FJ centre bounds \eqref{eq:wmt:FJ-asym}--\eqref{eq:wmt:FJ-normal} & $n,R,\rstar$ \\
$C_3(n,R,\rstar)$ & oriented radial trace \eqref{eq:wmt:radial-trace} & $n,R,\rstar$ \\
$C_4(n,R,\rstar)$ & good-level homothetic trace \eqref{eq:wmt:htrace-good} & $n,R,\rstar$ \\
$C_8(n,R,\rstar)$ & CS defect \eqref{eq:wmt:CS-defect} & $n,R,\rstar$ \\
$L_0(n,R,\rstar)$ & uniform metric Lipschitz \eqref{eq:wmt:A0} & $n,R,\rstar$ \\
$K_1(n,\rstar)=2/(\omega_n\rstar^n)$ & integrated Talenti gap \eqref{eq:wmt:Talenti-gap} & $n,\rstar$ \\
$K_2(n,\rstar)=4n/(\omega_n\rstar^{n+1})$ & bad-$(-t')$ measure \eqref{eq:wmt:Btau0-bound} & $n,\rstar$ \\
$\tau_0=\rstar/(2n)$ & bad-$(-t')$ threshold \eqref{eq:wmt:def-Btau0} & $n,\rstar$ \\
$c_0(n,R,\rstar)=4n/(\omega_n\rstar^{n+1})$ & window-scale constant, Thm.~\ref{thm:wmt:trace} & $n,R,\rstar$ \\
$C_*^{\mathrm{act}}(n,R,\rstar)$ & integrated action \eqref{eq:wmt:52int} & $n,R,\rstar$ \\
$C_*^{\mathrm{kin}}(n,R,\rstar)$ & Lebesgue kinetic \eqref{eq:wmt:K} & $n,R,\rstar$ \\
$C_*(n,R,\rstar)$ & weighted metric trace \eqref{eq:wmt:Cstar} & $n,R,\rstar$ \\
$C_{\mathrm{lay}}(n,R,\rstar)$ & boundary-layer volume, Rem.~\ref{rem:wmt:layer-volume} & $n,R,\rstar$ \\
\hline
\end{tabular}
\end{center}

All constants are polynomial in $(1/\rstar, R, n)$. Since
$\kappa=\kappa(n,R,\rstar)$ has been pinned at the start of \S4.4 to
the canonical value of the notation register \S12 (here we use
$\kappa=1$), every constant in the table has signature
$(n,R,\rstar)$; there is no residual $\kappa$-dependence. The output
$\rhad$ lies in the $O(\sqrt{\delta})$-window
\[
   \rhad\in[\rhodelta-c_0\sqrt{\delta},\rhodelta-c_0\sqrt{\delta}/2]
   \subset[\rhodelta-c_0\sqrt{\delta},\rhodelta],
\]
so $\rhodelta-\rhad\le c_0\sqrt{\delta}$ and (combined with
$1-\rhodelta=\kappa\sqrt{\delta}$)
$1-\rhad\le(\kappa+c_0)\sqrt{\delta}$. The boundary layer
$\Om\setminus E_\rhad$ has volume bounded by
$C_{\mathrm{lay}}(n,R,\rstar)\sqrt{\delta}$
(Remark~\ref{rem:wmt:layer-volume}), to be squared in the
boundary-layer transfer step
(Section~\ref{subsec:boundary-layer-transfer}). This delivers exactly
the source-map item E4.6 input required by Agent~13's
Theorem~\ref{thm:metric-trace-input}.
